% GENERATED FILE. Edit canonical lessons, then run npm run generate:books.
% canonical-source-hash: fnv1a64:3f3beff867518485
% canonical-lessons: JA-C420-tesuuryou, JA-C420-ryoushuusho, JA-C420-souryou, JA-C420-shiharau, JA-C420-haraimodosu

\chapter{Receipts and Fees}
\label{ch:ja-receipts-and-fees}
\hlchaptermodality{420}

\begin{chapteropening}
\textbf{By the end of this chapter:} \emph{I can pay, ask for a receipt and get a refund.}

Complete the last lesson of chapter 420: I can pay, ask for a receipt and get a refund.
\end{chapteropening}

\section[\texorpdfstring{\ja{てすうりょう} (tesūryō)}{てすうりょう (tesūryō)}]{\ja{てすうりょう} (tesūryō) — a handling fee}

\label{lesson:JA-C420-tesuuryou}



\begin{quote}
Before the new one: say the Japanese for one's full name, then the Japanese for to accept applications.
\end{quote}

\subsection*{You'll want to know: \ja{てすうりょう}}
\textbf{\ja{てすうりょう}} — \emph{tesūryō} — \textquotedblleft{}a handling fee\textquotedblright{}.

\begin{quote}
\textbf{\ja{ぎんこうで} \ja{てすうりょうを} \ja{はらいました。}} — \emph{ginkō de tesūryō o haraimashita} — \textquotedblleft{}I paid a handling fee at the bank.\textquotedblright{}
\end{quote}

\begin{cousinweb}
\textbf{\ja{てすうりょう}} (\emph{tesūryō}) is built from both layers of the vocabulary: the native Japanese (wago) \emph{te}, \textquotedblleft{}hand\textquotedblright{}, inherited rather than borrowed; the Sino-Japanese \emph{sū}, \textquotedblleft{}number\textquotedblright{} (Middle Chinese \emph{srjuH}), an on'yomi borrowed from Chinese; and the Sino-Japanese \emph{ryō}, \textquotedblleft{}fee\textquotedblright{} (\emph{lew}), an on'yomi borrowed from Chinese (JMdict; KANJIDIC2; Middle Chinese in Baxter's transcription).
\end{cousinweb}

\begin{grammarlens}[title={What you've built}]
One more word for getting around.
\end{grammarlens}

\subsection*{Your turn}
\begin{itemize}
\raggedright
  \item \emph{Say it:} \emph{tesūryō}
  \item \emph{Say it:} \emph{tesūryō}, once more
  \item \emph{Say it:} \emph{ginkō de tesūryō o haraimashita} — I paid a handling fee at the bank.
  \item \emph{Recall:} say the Japanese for one's full name, then the Japanese for to accept applications, then say \emph{tesūryō} again
\end{itemize}

\begin{tcolorbox}[breakable,colback=teal!4,colframe=teal!35!black,title={Before you move on}]
What does \ja{てすうりょう} mean? (\textquotedblleft{}A handling fee\textquotedblright{}.) Say it once more.
\end{tcolorbox}
\section[\texorpdfstring{\ja{りょうしゅうしょ} (ryōshūsho)}{りょうしゅうしょ (ryōshūsho)}]{\ja{りょうしゅうしょ} (ryōshūsho) — a receipt}

\label{lesson:JA-C420-ryoushuusho}



\begin{quote}
Before the new one: say the Japanese for to accept applications, then the Japanese for a handling fee.
\end{quote}

\subsection*{You'll want to know: \ja{りょうしゅうしょ}}
\textbf{\ja{りょうしゅうしょ}} — \emph{ryōshūsho} — \textquotedblleft{}a receipt\textquotedblright{}.

\begin{quote}
\textbf{\ja{りょうしゅうしょを} \ja{ください。}} — \emph{ryōshūsho o kudasai} — \textquotedblleft{}A receipt, please.\textquotedblright{}
\end{quote}

\begin{cousinweb}
\textbf{\ja{りょうしゅうしょ}} (\emph{ryōshūsho}) is Sino-Japanese (kango): three Chinese characters read with their on'yomi, readings borrowed from Chinese. They are \emph{ryō}, \textquotedblleft{}jurisdiction\textquotedblright{} (Middle Chinese \emph{ljengX}), \emph{shū}, \textquotedblleft{}obtain\textquotedblright{} (\emph{syuw}) and \emph{sho}, \textquotedblleft{}write\textquotedblright{} (\emph{syo}) (KANJIDIC2; Middle Chinese in Baxter's transcription).
\end{cousinweb}

\begin{grammarlens}[title={What you've built}]
One more word for getting around.
\end{grammarlens}

\subsection*{Your turn}
\begin{itemize}
\raggedright
  \item \emph{Say it:} \emph{ryōshūsho}
  \item \emph{Say it:} \emph{ryōshūsho}, once more
  \item \emph{Say it:} \emph{ryōshūsho o kudasai} — A receipt, please.
  \item \emph{Recall:} say the Japanese for to accept applications, then the Japanese for a handling fee, then say \emph{ryōshūsho} again
\end{itemize}

\begin{tcolorbox}[breakable,colback=teal!4,colframe=teal!35!black,title={Before you move on}]
What does \ja{りょうしゅうしょ} mean? (\textquotedblleft{}A receipt\textquotedblright{}.) Say it once more.
\end{tcolorbox}
\section[\texorpdfstring{\ja{そうりょう} (sōryō)}{そうりょう (sōryō)}]{\ja{そうりょう} (sōryō) — postage; a delivery charge}

\label{lesson:JA-C420-souryou}



\begin{quote}
Before the new one: say the Japanese for a handling fee, then the Japanese for a receipt.
\end{quote}

\subsection*{You'll want to know: \ja{そうりょう}}
\textbf{\ja{そうりょう}} — \emph{sōryō} — \textquotedblleft{}postage; a delivery charge\textquotedblright{}.

\begin{quote}
\textbf{\ja{おみやげの} \ja{そうりょうは} \ja{いくらですか。}} — \emph{omiyage no sōryō wa ikura desu ka} — \textquotedblleft{}How much is the postage for the souvenirs?\textquotedblright{}
\end{quote}

\begin{cousinweb}
\textbf{\ja{そうりょう}} (\emph{sōryō}) is Sino-Japanese (kango): two Chinese characters read with their on'yomi, readings borrowed from Chinese. They are \emph{sō}, \textquotedblleft{}send\textquotedblright{} (Middle Chinese \emph{suwngH}) and \emph{ryō}, \textquotedblleft{}fee\textquotedblright{} (\emph{lew}) (KANJIDIC2; Middle Chinese in Baxter's transcription).
\end{cousinweb}

\begin{grammarlens}[title={What you've built}]
One more word for getting around.
\end{grammarlens}

\subsection*{Your turn}
\begin{itemize}
\raggedright
  \item \emph{Say it:} \emph{sōryō}
  \item \emph{Say it:} \emph{sōryō}, once more
  \item \emph{Say it:} \emph{omiyage no sōryō wa ikura desu ka} — How much is the postage for the souvenirs?
  \item \emph{Recall:} say the Japanese for a handling fee, then the Japanese for a receipt, then say \emph{sōryō} again
\end{itemize}

\begin{tcolorbox}[breakable,colback=teal!4,colframe=teal!35!black,title={Before you move on}]
What does \ja{そうりょう} mean? (\textquotedblleft{}Postage; a delivery charge\textquotedblright{}.) Say it once more.
\end{tcolorbox}
\section[\texorpdfstring{\ja{しはらう} (shiharau)}{しはらう (shiharau)}]{\ja{しはらう} (shiharau) — to pay a bill; to pay a charge}

\label{lesson:JA-C420-shiharau}



\begin{quote}
Before the new one: say the Japanese for a receipt, then the Japanese for postage.
\end{quote}

\subsection*{You'll want to know: \ja{しはらう}}
\textbf{\ja{しはらう}} — \emph{shiharau} — \textquotedblleft{}to pay a bill; to pay a charge\textquotedblright{}.

\begin{quote}
\textbf{\ja{げんきんで} \ja{しはらいました。}} — \emph{genkin de shiharaimashita} — \textquotedblleft{}I paid in cash.\textquotedblright{}
\end{quote}

Polite \textbf{\ja{しはらいます}}, \emph{shiharaimasu}; a godan verb. Past (polite): \textbf{\ja{かいしゃが} \ja{りょひを} \ja{しはらいました。}}, \emph{kaisha ga ryohi o shiharaimashita}.

\begin{cousinweb}
\textbf{\ja{しはらう}} (\emph{shiharau}) is built from both layers of the vocabulary: the Sino-Japanese \emph{shi}, \textquotedblleft{}support\textquotedblright{} (Middle Chinese \emph{tsye}), an on'yomi borrowed from Chinese, and the native Japanese (wago) \emph{harau}, \textquotedblleft{}to pay\textquotedblright{}, inherited rather than borrowed (JMdict; KANJIDIC2; Middle Chinese in Baxter's transcription).
\end{cousinweb}

\begin{grammarlens}[title={What you've built}]
One more word for getting around.
\end{grammarlens}

\subsection*{Your turn}
\begin{itemize}
\raggedright
  \item \emph{Say it:} \emph{shiharau}
  \item \emph{Say it:} \emph{shiharau}, once more
  \item \emph{Say it:} \emph{genkin de shiharaimashita} — I paid in cash.
  \item \emph{Recall:} say the Japanese for a receipt, then the Japanese for postage, then say \emph{shiharau} again
\end{itemize}

\begin{tcolorbox}[breakable,colback=teal!4,colframe=teal!35!black,title={Before you move on}]
What does \ja{しはらう} mean? (\textquotedblleft{}To pay a bill; to pay a charge\textquotedblright{}.) Say it once more.
\end{tcolorbox}
\section[\texorpdfstring{\ja{はらいもどす} (haraimodosu)}{はらいもどす (haraimodosu)}]{\ja{はらいもどす} (haraimodosu) — to refund}

\label{lesson:JA-C420-haraimodosu}



\begin{quote}
Before the new one: say the Japanese for postage, then the Japanese for to pay a bill.
\end{quote}

\subsection*{You'll want to know: \ja{はらいもどす}}
\textbf{\ja{はらいもどす}} — \emph{haraimodosu} — \textquotedblleft{}to refund\textquotedblright{}.

\begin{quote}
\textbf{\ja{うんきゅうに} \ja{なったので、えきが} \ja{うんちんを} \ja{はらいもどします。}} — \emph{unkyū ni natta node, eki ga unchin o haraimodoshimasu} — \textquotedblleft{}The service was cancelled, so the station is refunding the fare.\textquotedblright{}
\end{quote}

Polite \textbf{\ja{はらいもどします}}, \emph{haraimodoshimasu}; a godan verb. Past (polite): \textbf{\ja{えきが} \ja{うんちんを} \ja{はらいもどしました。}}, \emph{eki ga unchin o haraimodoshimashita}.

\begin{cousinweb}
\textbf{\ja{はらいもどす}} (\emph{haraimodosu}) is a native Japanese word (wago), inherited rather than borrowed from Chinese, built from \emph{harai}, \textquotedblleft{}paying\textquotedblright{} (from \emph{harau}) and \emph{modosu}, \textquotedblleft{}to give back\textquotedblright{}; its kanji were attached later as kun'yomi (JMdict).
\end{cousinweb}

\begin{grammarlens}[title={What you've built}]
One more word for getting around.
\end{grammarlens}

\subsection*{Your turn}
\begin{itemize}
\raggedright
  \item \emph{Say it:} \emph{haraimodosu}
  \item \emph{Say it:} \emph{haraimodosu}, once more
  \item \emph{Say it:} \emph{unkyū ni natta node, eki ga unchin o haraimodoshimasu} — The service was cancelled, so the station is refunding the fare.
  \item \emph{Recall:} say the Japanese for postage, then the Japanese for to pay a bill, then say \emph{haraimodosu} again
\end{itemize}

\begin{tcolorbox}[breakable,colback=teal!4,colframe=teal!35!black,title={Before you move on}]
What does \ja{はらいもどす} mean? (\textquotedblleft{}To refund\textquotedblright{}.) Say it once more.
\end{tcolorbox}
